\nocite{0e533eef}
\nocite{3e24354b}
\nocite{37e40023}
\nocite{5cb8b651}
\nocite{9104a30f}
\nocite{99aba7f6}
\nocite{9fc084b2}
\nocite{c9fb3de7}
\nocite{ce883d84}
\nocite{d02f3620}
\nocite{eda79af1}
\nocite{f94459b0}
\nocite{5914f1ad}
In this section, we want to give the reader an idea of what a topological insulator is.
As the name suggests, both topology and insulators play a role; therefore, the section is essentially divided into two parts: a first part concerning ordinary insulators, and a second part concerning their topological aspects.
In the first part, we treat ordinary insulators.
We will present the standard mathematical model of insulators, leveraging a Bloch-Floquet transformation.
Explicitly using a Bloch-Floquet transformation is advantageous because it makes the topological information of the insulator accessible by allowing us to construct a vector bundle that exhibits a potentially interesting topology.
In fact, we will sketch the bundle construction after discussing insulators, as it provides the foundation for our treatment of topological insulators.
Equipped with this bundle, we move on to the second part.
There, after briefly discussing the phenomenon of topological phases and their K-theoretic classification, we examine the special case of so-called Chern insulators by analyzing their topological structure using topological K-theory on that bundle.
In particular, we will identify the integer quantum Hall effect as a characteristic of two-dimensional Chern insulators with the help of Chern classes.
\\
\begin{rem}
\label{rem:blochfloquet}
As a side note, the curious reader interested in the details of the aforementioned bundle construction via the Bloch-Floquet transform is referred to \cite{9fc084b2}, \cite{d02f3620}, and \cite{ce883d84}.
The latter paper discusses the generalization of the Bloch-Floquet transformation for arbitrary separable Hilbert spaces.
\\
\phantom{proven}
\hfill
$\square$
\\
\end{rem}
Let us begin by discussing ordinary insulators in this first part.
In physics, an insulator is a crystalline solid with very low electrical conductivity.
To make this definition accessible to theoretical considerations, we will now present a mathematical model of a {\glqq}crystalline solid{\grqq} and explain what it means for such a solid to exhibit {\glqq}low electrical conductivity.{\grqq}
\\
For many purposes, it suffices to consider a crystalline solid as an (infinite) periodic arrangement of atomic nuclei.
The standard approach models this periodicity using a so-called lattice and the arrangement of atomic nuclei as a so-called basis\footnote{This is a rather unfortunate naming convention in the context of mathematical models involving vector spaces, which use the term {\glqq}basis{\grqq} differently.}, assigning the type of atomic nucleus to a position within an area modulo the periodicity of the lattice.
This model is usually justified if the number of atomic nuclei is extremely large and the temperature is near (absolute) zero.
While the former is almost always true, the latter typically only occurs in highly specialized environments; however, for an introductory approach to (topological) insulators, this assumption is acceptable.
\\
Formally, an \textbf{($n$-dimensional) lattice} is a subspace $L$ of $\mathbb{R}^{n}$ such that there is a basis $l_{1},\ldots,l_{n} \in \mathbb{R}^{n}$ of $\mathbb{R}^{n}$ satisfying $L = \mathrm{span}_{\mathbb{Z}}(l_{1},\ldots,l_{n})$. In particular, $L$ is a group operating on $\mathbb{R}^{n}$ by $(l,x) \mapsto l + x$ and $F_{L} := \mathrm{span}_{[0,1)}(l_{1},\ldots,l_{n})$ is a fundamental domain w.r.t. that group action\footnote{Recall that a fundamental domain w.r.t. a group action $\phi$ is a connected subspace of the space on which $\phi$ operates, containing precisely one element from each orbit of $\phi$.} (representing an area modulo the lattice's periodicity).
Next, let $\mathfrak{C}$ (representing the positions of atomic nuclei) be a finite subset of $F_{L}$, and let $\mathsf{A}$ (representing the types of atomic nuclei) be a finite set. A \textbf{c-basis (for $L$)} is then defined as a function $\mathfrak{C} \to \mathsf{A}$.
A \textbf{crystal structure} is thus a lattice $L$ paired with a c-basis for $L$.
At this point, we need to introduce an important concept regarding crystal structures: the Brillouin zone (and its close relationship to momentum space).
To this end, we note that for any given lattice $L$, there exists a\footnote{Recall that $\mathbb{R}^{n}$ is self-dual; that is, $(\mathbb{R}^{n})^{\prime}= \mathcal{L}(\mathbb{R}^{n},\mathbb{R}) = \mathbb{R}^{n}$ as vector spaces.} \textbf{dual lattice $L^{\prime}$} defined by
\begin{align*}
  L^{\prime}
  &=
  \lbrace
      k
      \in
      \mathbb{R}^{n}
    \colon
      \forall
      x
      \in
      L
      \,
      (k \vert x)
      \in 
      2\pi
      \mathbb{Z}
  \rbrace
\end{align*}
In particular, there is a dual basis $l_{1}^{\prime},\ldots,l_{n}^{\prime} \in \mathbb{R}^{n}$ defined by $(l_{i}^{\prime} \vert l_{j}) = 2\pi\delta_{ij}$.
The \textbf{Brillouin zone $B_{L}$} is the topological closure of the fundamental domain $F_{{L}^{\prime}} = \mathrm{span}_{[-1/2,1/2)}(l_{1}^{\prime},\ldots,l_{n}^{\prime})$, written $B_{L} = \mathrm{cl}(F_{L^{\prime}})$.
\\
Having established a model for crystalline solids, we can now take the quantum mechanical nature of the atomic nuclei into account.
In our model, the atomic nuclei and their periodic occurrences generate a periodic potential $V$, such as a superposition of Coulomb-like potentials.
Moreover, it is clear that the quantum system is subject to at least a translational symmetry due to this periodicity.
\\
Formally, recall that a \textbf{symmetry (transformation)} on a complex and separable Hilbert space $(\mathcal{H},(\cdot \vert \cdot)_{\mathcal{H}})$ is a surjective function $S \colon \mathcal{H} \to \mathcal{H}$ such that
\begin{align*}
  \vert
    (S(\psi_{1}) \vert S(\psi_{2}))
  \vert
  &=
  \vert
    (\psi_{1} \vert \psi_{2})
  \vert
\end{align*}
for all $\psi_{1},\psi_{2} \in \mathcal{H}$ and that these symmetry transformations together with composition form a group: the \textbf{symmetry group $G_{\mathcal{H}}$}.
According to Wigner's Theorem, there is a unitary or antiunitary operator $U_{S} \colon \mathcal{H} \to \mathcal{H}$ for each $S \in G_{\mathcal{H}}$ such that
\begin{align*}
  S(\psi)
  &=
  \exp(\mathrm{i}\varphi) U_{S}(\psi)
\end{align*}
for all $\psi \in \mathcal{H}$ and some $\exp(\mathrm{i}\varphi) \in \mathbb{C}$.
Further recall that a quantum system is represented by a self-adjoint, linear operator $H \colon \mathrm{dom}(H) \to \mathcal{H}$ - the Hamilton operator - where $\mathrm{dom}(H)$ is a dense subset of the complex and separable Hilbert space $\mathcal{H}$.
A \textbf{symmetry (of $H$)} is then an $S \in G_{\mathcal{H}}$ such that
\begin{align*}
  U_{S}(\mathrm{dom}(H))
  &=
  \mathrm{dom}(H)
\end{align*}
and such that $U_{S}$ either commutes or anticommutes with $H$, depending on whether $U_{S}$ is unitary or antiunitary.
Restricting our attention to $\mathcal{H} = L_{2}(\mathbb{R}^{n},\mathbb{C})$, we can consider \textbf{periodic potentials} - i.e., $2$-locally Lebesgue integrable functions $V \in L_{2,\textrm{loc}}(\mathbb{R}^{n},\mathbb{R})$ such that $V(x + l) = V(x)$ for all $l \in L$ and some $n$-dimensional lattice $L$.
Then the Hamilton operator we are interested in,
\begin{align*}
  H
  &:=
  H_{0}
  +
  M_{V}
  :=
  -
  \frac{1}{2}
  \Delta
  +
  M_{V}
\end{align*}
is a self-adjoint operator, provided $M_{V}$ is the multiplication operator associated with $V$ and $\Delta$ is the Laplace operator.
Since $H$ commutes with the translation operators defined by $T_{l}\psi := \psi(\cdot - l)$, it possesses at least a translational symmetry.
\\
To fully understand crystalline solids quantum mechanically, we must solve the Schr{\"o}dinger equation for these Hamiltonians.
We present a method leveraging the aforementioned Bloch-Floquet transformation alongside spectral theory.
For Schwartz functions $f \in \mathcal{S}(\mathbb{R}^{n},\mathbb{C})$, one defines the \textbf{modified Bloch-Floquet transformation}
\begin{align*}
  \mathcal{U}
  \colon
  \mathcal{S}(\mathbb{R}^{n},\mathbb{C})
  \to
  L_{2}(B_{L},L_{2}(\mathrm{cl}(F_{L}),\mathbb{C}))
  ,\qquad
  f
  \mapsto
  (\mathcal{U}(f))(k)
  &:=
  \frac{1}{\sqrt{\mathrm{vol}_{n}(B_{L})}}
  \sum_{l \in L}
  \exp(\mathrm{i}(k \vert \cdot - l))
  f(\cdot - l)
\end{align*}
$\mathcal{U}$ can be extended to a unitary, linear operator on $L_{2}(\mathbb{R}^{n},\mathbb{C})$.
The modified Bloch-Floquet transformation plays a role for $H$ similar to what the Fourier transform does for $-\frac{1}{2}\Delta$; it helps solve the Schr{\"o}dinger equation for $H$ by mapping the problem into momentum space.
Specifically, $\mathcal{U}H\mathcal{U}^{-1}$ equals the \textit{fibered operator} $\int_{k \in B_{L}}^{\oplus}\tilde{H}(k)$, where $\tilde{H}$ is a specific function on $B_{L}$.
The notation $\int_{k \in B_{L}}^{\oplus}$ denotes a \textit{direct integral}.
Without delving into a formal definition, it suffices to say that direct integrals are a continuous generalization of the direct sum.
However, the important thing to note is that
\begin{align*}
  \tilde{H}(k)
  &=
  -
  \mathrm{i}(\nabla + k \cdot \mathrm{id})^{2}/2
  +
  M_{V}
\end{align*}
is a densely defined, self-adjoint, linear operator for each (crystal momentum) $k \in B_{L}$, with its domain being the corresponding second Sobolev space:
\begin{align*}
  \mathrm{dom}(\tilde{H}(k))
  &=
  H_{2}(\mathrm{cl}(F_{L}),\mathbb{C})
\end{align*}
Moreover, by spectral theory, the (energy) spectrum $\sigma(\tilde{H}(k_{0}))$ of $\tilde{H}(k_{0})$ for each $k_{0} \in B_{L}$ is a pure point spectrum; i.e., it consists solely of the eigenvalues $E_{i}(k_{0}) \in \mathbb{R}$ (where $i \in \mathbb{N}$) of $\tilde{H}(k_{0})$.
Without loss of generality, we assume $E_{i}(k_{0}) \leq E_{i+1}(k_{0})$ for all $i$.
Consequently,
\begin{align*}
  \mathrm{im}(E_{i})
  &=
  \bigcup_{k \in B_{L}}\lbrace E_{i}(k) \rbrace
\end{align*}
is a compact interval for each $i$, since it can be shown that $E_{i}$ is a continuous function on $B_{L}$.
The image $\mathrm{im}(E_{i})$ is called the \textbf{($i$-th spectral) band}, and the (countable) union of all these bands ultimately equals the full spectrum $\sigma(H)$ of $H$.
If two adjacent bands - that is, bands with indices $i$ and $i+1$ - have an empty intersection, one says there is a \textbf{(spectral) gap} between those bands.
All in all, this yields the so-called band model of crystalline solids, which is well-known in physics.
\\
We make two remarks at this point:
\begin{enumerate}
\item
  Note that the \textit{direct integral decomposition} allows us to decompose the spectrum of $H$ according to the \textit{crystal momentum $k$}.
This will make certain topological information apparent in the second part.
\item
  The ideas presented above also apply to Hamilton operators where $H_{0}$ is replaced by $\frac{1}{2}(-\mathrm{i}\nabla + M_{A})^{2}$ or $\frac{1}{2}((-\mathrm{i}\nabla + M_{A}) \cdot \sigma)^{2}$, provided $A \colon \mathbb{R}^{n} \to \mathbb{R}^{n}$ satisfies $A(x + l) = A(x)$ for $l \in L$, and $\sigma = (\sigma_{1},\sigma_{2},\sigma_{3})$ represents the standard Pauli matrices.
Here, $A$ abstracts a magnetic vector potential, and $\sigma$ accounts for the particle's spin.
The core concepts of the modified Bloch-Floquet transform and spectral bands remain unchanged; the modifications are minor.
%\item
%  Using the tensor product $\otimes$ we can describe $N$ non-interacting, identical particles by the operator $\sum_{i=1}^{N}H_{(i)}$ defined on $\bigotimes_{i=1}^{N}\mathrm{dom}(H_{(i)})$ where $H_{(i)} = \bigotimes_{j=1}^{N}O_{j}$ with $O_{i} = H$ and the other $O_{j}$ being the identity.
\end{enumerate}
Equipped with the concept of energy bands, we can define what constitutes an insulator: if the \textit{Fermi level} lies within a band gap, we say the crystal structure is an \textbf{insulator}.
We refrain from giving a formal definition of the Fermi level here. Instead, we content ourselves with the heuristic explanation that at absolute zero, the Fermi level is the energy threshold where all states below it are fully occupied, and all states above it are entirely unoccupied.
This justifies our definition of an insulator - at absolute zero, at least - because electrical conductivity requires unoccupied states within a band to facilitate charge transport. In a fully occupied band, for every electron in state $k$, there is a corresponding electron in state $-k$, resulting in a net electron momentum of $0$ for that band.
Consequently, no net electron transport occurs, and therefore no electric current can flow.
\\
%\\
%For the sake of completeness we will give as somewhat informal definition of the Fermi-level.
%The definition will make use of the finiteness of electrons in a crystalline solid.
%This is not reflected in our model so far as we implicitly made the (correct) assumption that the number of atomic nuclei is very large and hence approximated the large number with infinity.
%Most of the time the approximation of crystalline solids with crystal structures exhibiting a band structure on a single particle level actually suffices our needs.
%Indeed, when nothing other is said, this is assumed.
%But for our Fermi-level definition we need something slightly more realistic in that regard.
%Intersecting the $n$-dimensional lattice with an $n$-dimensional cube will do because the number of atomic nuclei is then finite and hence the number of electrons brought in by them to make the crystalline solid electrically neutral must be finite, too.\footnote{in addition, it models boundary as well as a finite volume} We assume these $N$ electrons as non-interacting fermions feeling the potential generated by the (finite number of) atomic nuclei in the crystalline solid.
%Hence the \textbf{(total) chemical potential} or \textbf{Fermi-level} $E_{F}$ of the this system of fermions is a well-defined thermodynamic quantity: the derivative of the \textit{internal energy} with respect to the particle (or more specifically electron) number $N$.
\\
Finally, we arrive at the promised discussion of vector bundles.
Decomposing the spectrum of $H$ according to the crystal momentum allows us to construct a vector bundle.
To this end, let $\tilde{I} \subset \mathbb{R}$ be an interval.
Furthermore, for $I := \tilde{I} \cap \mathbb{N}$, we define
\begin{align*}
  \sigma_{I}^{k_{0}}
  :=
  \bigcup_{i \in I}
  \lbrace
    E_{i}(k_{0})
  \rbrace
  \qquad
  &\text{and}
  \qquad
  \sigma_{I}
  :=
  \bigcup_{k \in B_{L}}
  \sigma_{I}^{k}
  =
  \bigcup_{i \in I}
  \mathrm{im}(E_{i})
\end{align*}
Using functional calculus, we obtain the spectral projections
\begin{align*}
  \tilde{P}^{I}(k)
  &:=
  \chi_{\sigma_{I}^{k}}(\tilde{H}(k))
\end{align*}
where $\chi$ denotes the characteristic function.
If the cardinality of $I$ is finite, and there is a spectral gap above the $\max(I)$-th band and below the $\min(I)$-th band, the family $\lbrace \tilde{P}^{I}(k) \colon k \in B_{L} \rbrace$ and the overall spectral projection
\begin{align*}
  P^{I}
  &:=
  \chi_{\sigma_{I}}(H)
\end{align*}
respectively, can be identified with a complex vector bundle over $\mathbb{R}^{n}/L^{\prime}$ - which is topologically the $n$-torus $\mathbb{T}^{n}$.
To construct the bundle, we define an equivalence relation on $\mathbb{R}^{n} \times L_{2}(\mathbb{R}^{n},\mathbb{C})$ by
\begin{align*}
  (k,\psi)
  \sim
  (\hat{k},\hat{\psi})
  \qquad
  &:\Leftrightarrow
  \qquad
  \exists
  l^{\prime}
  \in
  L^{\prime}
  \colon
  \,
  \hat{k}
  =
  k
  -
  l^{\prime}
  \quad
  \land
  \quad
  \hat{\psi}(\cdot)
  =
  \exp(-\mathrm{i}(l^{\prime} \vert \cdot))
  \psi(\cdot)
\end{align*}
The desired total space is then given by
\begin{align*}
  E_{P^{I}}
  &:=
  \left\lbrace
      [(k,\phi)]
      \in
      \left(
        \mathbb{R}^{n}
        \times
        L_{2}(\mathbb{R}^{n},\mathbb{C})
      \right)
      /
      \sim
    \colon
      \phi
      \in
      \mathrm{Ran}(\tilde{P}^{I}(k))
  \right\rbrace
\end{align*}
When studying the conductivity properties of crystalline solids and their associated topology, it is natural to examine the vector bundle $E_{P^{F}}$.
Here, $F$ is the set of band indices corresponding to energies below the Fermi level, since, as we have seen, the electronic properties are governed by the bands near the Fermi level.
In fact, this bundle is crucial for understanding topological phases, as we will see in the second part.
\\\\
In the second part, we want to talk about topological insulators.
Heuristically, a topological insulator is a material that is insulating in its bulk but can feature a {\glqq}conducting{\grqq} surface in certain {\glqq}topological phases{\grqq}.
These phases are characterized by so-called symmetry-protected {\glqq}surface states.{\grqq}
This means that surface states from different topological phases cannot be smoothly {\glqq}deformed{\grqq} into one another without breaking the underlying symmetries of the insulator.
One of these phases is designated as {\glqq}topologically trivial{\grqq} and corresponds to an ordinary insulator, whereas the {\glqq}non-trivial{\grqq} phases correspond to insulators exhibiting these conducting surface states.
\\
With all these quotation marks, this is clearly far from a rigorous definition.
In fact, providing a mathematically rigorous definition is difficult.
To complicate matters, our earlier definition of an insulator is merely a first-order approximation.
For example, if the temperature is significantly above absolute zero, the lattice is far from perfect, and our description may break down.
There are also other physical quirks that can render this simplified model invalid.
\cite{5cb8b651} and \cite{f94459b0} address these problems in detail by applying methods from non-commutative geometry.
This approach yields a much more robust model and represents the modern perspective on topological insulators.
Nonetheless, we will aim to understand topological insulators through the lens of so-called Chern insulators, utilizing more classical theories: the insulator model described in the first part, K-theory, and Chern classes.
\\
The paper \cite{c9fb3de7} develops a classification of the topological phases for \textit{gapped systems of free fermions subjected to symmetries in $A \subset G_{\mathcal{H}}$} using operator K-theory.
Operator K-theory is a generalization of topological K-theory where so-called \textit{Banach algebras} replace the topological spaces.
Depending on the symmetries of the quantum mechanical system, one can construct appropriate algebras and apply operator K-theory to derive associated groups.
The group elements then represent the different topological phases of the gapped systems with the specified symmetries in $A$.
Now, insulators as defined here are precisely such gapped systems of free fermions.
For $n$-dimensional insulators subjected to time-reversal symmetry, charge conjugation symmetry, a combination of both, or neither, the topological phase classification is summarized in Table \ref{tab:pt}.
\begin{table}[h!]
\begin{tabular}{cccccccccc}
  \hline
  $m$
  &
  Cartan label
  &
  $n=0$
  &
  $n=1$
  &
  $n=2$
  &
  $n=3$
  &
  $n=4$
  &
  $\cdots$
  &
  $n=8$
  &
  $\cdots$
  \\
  \hline
  $0$
  &
  AI
  &
  $\mathbb{Z}$
  &
  $0$
  &
  $0$
  &
  $0$
  &
  $\mathbb{Z}$
  &
  &
  $\mathbb{Z}$
  &
  \\
  $1$
  &
  BDI
  &
  $\mathbb{Z}_{2}$
  &
  $\mathbb{Z}$
  &
  $0$
  &
  $0$
  &
  $0$
  &
  &
  $\mathbb{Z}_{2}$
  &
  \\
  $2$
  &
  D
  &
  $\mathbb{Z}_{2}$
  &
  $\mathbb{Z}_{2}$
  &
  $\mathbb{Z}$
  &
  $0$
  &
  $0$
  &
  &
  $\mathbb{Z}_{2}$
  &
  \\
  $3$
  &
  DIII
  &
  $0$
  &
  $\mathbb{Z}_{2}$
  &
  $\mathbb{Z}_{2}$
  &
  $\mathbb{Z}$
  &
  $0$
  &
  &
  $0$
  &
  \\
  $4$
  &
  AII
  &
  $\mathbb{Z}$
  &
  $0$
  &
  $\mathbb{Z}_{2}$
  &
  $\mathbb{Z}_{2}$
  &
  $\mathbb{Z}$
  &
  &
  $\mathbb{Z}$
  &
  \\
  $5$
  &
  CII
  &
  $0$
  &
  $\mathbb{Z}$
  &
  $0$
  &
  $\mathbb{Z}_{2}$
  &
  $\mathbb{Z}_{2}$
  &
  &
  $0$
  &
  \\
  $6$
  &
  C
  &
  $0$
  &
  $0$
  &
  $\mathbb{Z}$
  &
  $0$
  &
  $\mathbb{Z}_{2}$
  &
  &
  $0$
  &
  \\
  $7$
  &
  CI
  &
  $0$
  &
  $0$
  &
  $0$
  &
  $\mathbb{Z}$
  &
  $0$
  &
  &
  $0$
  &
  \\
  \hline
  \hline
  $0$
  &
  A
  &
  $\mathbb{Z}$
  &
  $0$
  &
  $\mathbb{Z}$
  &
  $0$
  &
  $\mathbb{Z}$
  &
  &
  $\mathbb{Z}$
  &
  \\
  $1$
  &
  AIII
  &
  $0$
  &
  $\mathbb{Z}$
  &
  $0$
  &
  $\mathbb{Z}$
  &
  $0$
  &
  &
  $0$
  \\
  \hline
  \\
\end{tabular}
\caption{Periodic table for topological phases of insulators in $n$ dimension with symmetry data indicated by the Cartan label, see \cite{9104a30f} for example}
\label{tab:pt}
\end{table}
The groups appearing in Table \ref{tab:pt} are precisely $K_{\mathbb{K}}^{n-m}(\lbrace x_{0} \rbrace)$, which represents the topological K-theory of a point, where $\mathbb{K}$ equals $\mathbb{C}$ for Cartan labels A and AIII.
Otherwise, $\mathbb{K}$ equals $\mathbb{R}$.
\\
To gain a better intuition for this classification, we restrict our attention to insulators with Cartan label A, which lack both time-reversal and charge conjugation symmetries.
These are known as \textbf{Chern insulators} because their topological phases are classified by \textit{Chern numbers}.
The remainder of this section offers a brief glimpse into how this works.
\\
Recall that the topological phases of $n$-dimensional Chern insulators are classified by $K_{\mathbb{C}}^{n}(\lbrace x_{0} \rbrace)$.
On the other hand, for any such insulator, we have the vector bundle $E_{P^{F}}$ over the $n$-torus $\mathbb{T}^{n}$, which encapsulates its topological information.
Remarkably, the stable isomorphism class of $E_{P^{F}}$ completely determines its topological phase as a Chern insulator. As we saw in subsection \ref{sec:topktheory}, the relationship
\begin{align*}
  K_{\mathbb{C}}^{0}(\mathbb{T}^{n})
  &=
  K_{\mathbb{C}}^{n}(\lbrace x_{0} \rbrace)
  \oplus
  \bigoplus_{j=0}^{n-1}
  K_{\mathbb{C}}^{j}(\lbrace x_{0} \rbrace)^{\binom{n}{j}}
\end{align*}
holds.\footnote{Incidentally, the remaining summands relate to the phenomenon of \textit{weak} topological insulators.} This is a striking result, although we cannot yet relate it directly to geometry (as mentioned in subsection \ref{sec:topktheory}).
We will partially remedy this in the rest of the section.
\\
Specifically, if $X$ is a smooth manifold (like the torus $\mathbb{T}^{n}$), one can calculate the invariant using more analytical methods by adopting a differential geometric perspective on Chern classes.
While we will not dwell on it extensively, Chern classes do carry profound geometric meaning within the framework of topological K-theory.
However, this geometric aspect is most apparent in their differential geometric definition, where they are derived directly from the curvature of a bundle.
More precisely, if $(E,X,p,\mathbb{C}^{n})$ is a \textit{Hermitian}, \textit{smooth}, complex vector bundle over a compact manifold $X$, it possesses a \textit{curvature form} $\Omega_{E}$. One can then define the \textbf{$n$-th Chern class (of $E \in V_{\mathbb{C}}(X)$)} $c_{n}(E)$ for all $n \in \mathbb{N}$ via the characteristic polynomial
\begin{align*}
  \det
  \left(
    \frac{\mathrm{i}t\Omega_{E}}{2\pi}
    +
    \mathrm{id}
  \right)
  &=:
  \sum_{n=0}^{\infty}
  c_{n}(E)
  t^{n}
\end{align*}
For the first few terms, one obtains
\begin{align*}
  \sum_{n=0}^{\infty}
  c_{n}(E)
  t^{n}
  &=
  \mathrm{id}
  +
  \mathrm{i}
  \frac{\mathrm{tr}(\Omega_{E})}{2\pi}
  t
  +
  \frac{\mathrm{tr}(\Omega_{E}^{2}) - \mathrm{tr}(\Omega_{E})^{2}}{8\pi^{2}}
  t^{2}
  +
  \ldots
\end{align*}
where $\mathrm{tr}(\cdot)$ is the trace.
Furthermore, if the compact manifold $X$ is \textit{oriented} and has dimension $2d$, then for any integer partition $D := (d_{1},\ldots,d_{r}) \in \mathbb{N}^{r}$ of $d$ (meaning $d = \sum_{j}d_{j}$), the \textbf{$D$-th Chern number} is defined as $c_{D}(E) := c_{d_{1}}(E) \cdots c_{d_{r}}(E)$.
Importantly, one can show\footnote{This relates to homology: if you are curious, the relevant keywords are de Rham cohomology and Kronecker pairing.}
\begin{align*}
  \int_{X}
  c_{D}(E)
  \in
  \mathbb{Z}
\end{align*}
in this setting.
Finally, the topological phase of a Chern insulator in $n = 2d$ dimensions is determined by the $(d,0)$-th Chern number of the vector bundle $E_{P^{F}}$, which is calculated using the formula
\begin{align*}
  \mathrm{Ch}_{d}(P^{F})
  &:=
  \int_{\mathbb{T}^{n}}
  c_{1}(E_{P^{F}})^{d}
  \in
  K_{\mathbb{C}}^{n}(\lbrace x_{0} \rbrace)
\end{align*}
For $n = 2$, this simplifies to
\begin{align*}
  \mathrm{Ch}_{1}(P^{F})
  &=
  \frac{\mathrm{i}}{2\pi}
  \int_{\mathbb{T}^{2}}
  \mathrm{tr}
  \left(
    \tilde{P}^{F}
    \cdot
    [\partial_{1}\tilde{P}^{F},\partial_{2}\tilde{P}^{F}]
  \right)
  \mathrm{d}k
  \in
  K_{\mathbb{C}}^{2}(\lbrace x_{0} \rbrace)
\end{align*}
where $[\cdot,\cdot]$ is the commutator.
This models the integer quantum Hall effect, with the Hall conductance given by $\frac{e^{2}}{h}\mathrm{Ch}_{1}(P^{F})$.
Incidentally, these integrals - especially for $n = 2$ - bear a striking resemblance to the Gau{\ss}-Bonnet theorem: integrating a curvature-related quantity yields a topological invariant.
\\
We conclude with a remark regarding a link between K-theory and Chern classes.
\\
\begin{rem}
\label{rem:chchar}
There is a more general, axiomatic definition of Chern classes that applies to vector bundles over any compact Hausdorff base space.
It is formulated in terms of singular cohomology - a specific covariant functor $H_{\mathbb{Q}}^{2\ast}$ from $\mathbf{Top-cH}$ to $\mathbf{Ring}$. The $n$-th Chern class $c_{n}(E)$ of $E \in V_{\mathbb{C}}(X)$ is the $n$-th element in a unique sequence of functions
\begin{align*}
  \left(
    c_{n}
    \colon
    V_{\mathbb{C}}(X)
    \to
    H_{\mathbb{Q}}^{2\ast}(X)
  \right)_{n\in\mathbb{N}}
\end{align*}
satisfying certain axioms.
We will not elaborate on this axiomatization here, but we will take it for granted to define the Chern character.
To that end, let $s_{i}$ (for $i \in \mathbb{N}^{\times}$) denote the \textit{Newton polynomials}.
Then the \textit{Chern character $\mathrm{ch}(\cdot)$} is defined by
\begin{align*}
  \mathrm{ch}
  &\colon
   V_{\mathbb{C}}(X)
  \to
  H_{\mathbb{Q}}^{2\ast}(X)
  ,\qquad
  (E,X,p,\mathbb{C}^{n})
  \mapsto
  \mathrm{dim}(\mathbb{C}^{n})
  +
  \sum_{i=1}^{\infty}
  \frac{s_{i}(c_{1}(E),\ldots,c_{i}(E))}{i!}
\end{align*}
Here, the first few terms are given by
\begin{align*}
  \mathrm{ch}(E)
  &=
  n
  +
  c_{1}(E)
  +
  \frac{1}{2}
  (c_{1}(E)^{2} - 2c_{2}(E))
  +
  \ldots
\end{align*}
In particular, for a compact manifold $X$, we obtain
\begin{align*}
  \mathrm{ch}(E)
  &=
  \mathrm{tr}(\exp(\mathrm{i}\Omega_{E}/2\pi))
\end{align*}
as one can verify from the expression above for the first few terms.
In either case, the Chern character satisfies
\begin{align*}
  \mathrm{ch}(E_{1} \oplus E_{2})
  =
  \mathrm{ch}(E_{1})
  +
  \mathrm{ch}(E_{2})
  \qquad
  &\text{and}
  \qquad
  \mathrm{ch}(E_{1} \otimes E_{2})
  =
  \mathrm{ch}(E_{1})
  \mathrm{ch}(E_{2})
\end{align*}
Finally, one can define a natural transformation $\widetilde{\mathrm{ch}}$ from $K_{\mathbb{C}}^{0}$ to $H_{\mathbb{Q}}^{2\ast}$ satisfying
\begin{align*}
  \widetilde{\mathrm{ch}}([(E,\epsilon^{0})])
  &:=
  \mathrm{ch}(E)
\end{align*}
for $[(E,\epsilon^{0})] \in K_{\mathbb{C}}(X)$, which unveils the profound link between topological K-theory and Chern classes.
In the case of the integer quantum Hall effect, the Chern character of the corresponding bundle $E_{P^{F}}$ is
\begin{align*}
  \mathrm{ch}(E_{P^{F}})
  &=
  c_{1}(E_{P^{F}})
  +
  \mathrm{dim}(E_{P^{F}})
\end{align*}
reflecting the direct sum decomposition of $K_{\mathbb{C}}^{0}(\mathbb{T}^{2})$.
Note that the summand $K_{\mathbb{C}}^{0}(\lbrace x_{0} \rbrace)$ strictly determines the dimension of the representatives of $[E]_{s} \in K_{\mathbb{C}}^{0}(\mathbb{T}^{n})$.
\end{rem}
